% ============================================================
% Paper 1 — IEEE BIBM 2026
% B-JEPA Failure Analysis + GCV-Ridge + IV Framing
%
% Timeline:
%   April 19 2026  — arXiv preprint (establishes priority)
%   Jul–Aug  2026  — IEEE BIBM 2026 submission (primary target)
%   Oct      2026  — NeurIPS MLCB workshop (backup)
%   Dec      2026  — IEEE BIBM 2026 conference
%
% Target track: Topic 2.d (Gene Expression and Regulation)
%               Topic 5.a (Data Mining, ML, and AI)
%
% Target length: ~10 pages main + supplementary (IEEE BIBM format)
% ============================================================

\documentclass[11pt]{article}
\usepackage[margin=1in]{geometry}
\usepackage{amsmath, amssymb, amsthm}
\usepackage{graphicx}
\usepackage{hyperref}
\usepackage{booktabs}
\usepackage{natbib}
\usepackage{xcolor}

\newcommand{\todo}[1]{\textcolor{red}{[TODO: #1]}}

\title{Dense World Models Fail for Sparse Causal Inference:\\
       A Diagnostic Study of B-JEPA for Gene Regulatory Network Inference}
\author{Raja Giddi}
\date{April 2026}

\begin{document}
\maketitle

% ============================================================
\begin{abstract}
Self-supervised world models based on the Joint-Embedding Predictive Architecture
(JEPA) have shown strong performance on dense prediction tasks in vision.
We ask whether B-JEPA, a Bayesian adaptation of JEPA, can learn gene regulatory
network (GRN) structure from observational expression data without supervision.
It cannot --- and the failure is instructive.
We identify three separable, diagnosable failure modes:
(I) the ELBO collapses (KL $\approx 1.07$ across all epochs and architectures),
producing representations indistinguishable from prior samples;
(II) mean-pool context aggregation is architecturally mismatched to sparse local
causal inference, producing a weight initialisation $W_{\rm init}$ with no
per-TF-gene selectivity;
(III) mean-field variational inference collapses the horseshoe prior to ridge for
every edge, eliminating the intended sparsification.
A GCV-calibrated horseshoe ridge estimator --- computed analytically in seconds,
violating none of the three failure conditions --- outperforms GENIE3 on three of
four DREAM5 benchmark networks (AUPR: $+0.043$ on \textit{E.~coli},
$+0.035$ on \textit{S.~aureus} via ensemble, $+0.003$ on \textit{S.~cerevisiae}).
Ablations of BiJEPA, KAN Stage~2, and stochastic-masking JEPA confirm that
architectural modifications to B-JEPA do not resolve any of the three failures.
A direction-ensemble combining forward GCV, inverted masking, and B-JEPA achieves
the best overall AUPR (0.2615 on net1, beating GENIE3 by $+0.003$).
We conclude with a causal identification analysis framing DREAM5 experimental
conditions as natural instruments and specifying the architectural requirements
for a correct causal world model for GRN inference.
\end{abstract}

% ============================================================
\section{Introduction}

Gene regulatory network (GRN) inference is the problem of determining which
transcription factors (TFs) causally regulate which target genes from
observational or interventional expression data.
It is not a prediction problem and not a correlation problem: the objective is
to rank $D \times G$ candidate directed edges ($D$ TFs, $G$ genes) by causal
regulatory strength, evaluated by area under the precision-recall curve (AUPR).
The sparsity is extreme --- fewer than $0.2\%$ of candidate edges are true
regulatory interactions --- making GRN inference a sparse local causal
identification problem.

The Joint-Embedding Predictive Architecture (JEPA) \citep{lecun2022} proposes
that world models should predict in \emph{latent space} rather than pixel space,
using an EMA target encoder and a context predictor.
The natural hypothesis is that a self-supervised JEPA model, trained on
gene expression data, will learn representations that capture regulatory
co-variation --- effectively discovering GRN structure without edge-label
supervision.
B-JEPA \citep{huang2026} instantiates this hypothesis with a Bayesian Stage~2
horseshoe regression on the learned representations.

The hypothesis fails, and fails for three separable reasons.
This is not a failure of hyperparameter tuning or insufficient training --- we
show that the KL collapse, the dense aggregation mismatch, and the VI collapse
are structural properties of the architecture applied to this problem class.
Each failure points to a specific design principle that a correct causal world
model must satisfy.
Critically, a GCV-calibrated ridge estimator that satisfies all three design
principles analytically outperforms GENIE3 \citep{huynh2010} on three of four
DREAM5 benchmark networks \citep{marbach2012}.

This paper makes four contributions:
(1) a mechanistic diagnosis of three separable failure modes in B-JEPA for GRN
inference, with supporting quantitative evidence across all four DREAM5 networks;
(2) a GCV-calibrated analytical horseshoe ridge baseline that is faster and more
accurate than B-JEPA and GENIE3 on biological networks;
(3) an ensemble analysis showing that B-JEPA representations provide orthogonal
signal only in the well-conditioned regime ($n/D > 4$);
(4) a causal identification framework framing DREAM5 experimental conditions as
natural instruments and specifying the Causal B-JEPA architecture as the path
to calibrated credible intervals over regulatory edges.

% ============================================================
\section{Background}

\subsection{Gene Regulatory Network Inference and DREAM5}

The DREAM5 Network Inference Challenge \citep{marbach2012} provides four
benchmark networks with known gold-standard regulatory edges
(Table~\ref{tab:networks}).
The task is to rank all $D \times G$ TF--gene pairs by inferred regulatory
strength, evaluated by AUPR.
The $n/D$ ratio --- samples per TF predictor --- is the key regime indicator:
at $n/D < 2$, all linear methods are underdetermined and degrade; at $n/D > 4$,
regularised regression approaches the performance ceiling set by data quality.

\begin{table}[h]
\centering
\caption{DREAM5 network statistics.}
\label{tab:networks}
\begin{tabular}{lrrrrrr}
\toprule
Network & $n$ & $G$ & $D$ & $n/D$ & Pos.\ edges & Sparsity \\
\midrule
Net1 (\textit{in silico})     & 805 & 1{,}643 & 195 & 4.13 & 4{,}012 & 1.25\% \\
Net2 (\textit{S.~aureus})     & 160 & 2{,}810 &  99 & 1.62 &    428  & 0.15\% \\
Net3 (\textit{E.~coli})       & 805 & 4{,}511 & 334 & 2.41 & 2{,}066 & 0.14\% \\
Net4 (\textit{S.~cerevisiae}) & 536 & 5{,}950 & 333 & 1.61 & 3{,}940 & 0.20\% \\
\bottomrule
\end{tabular}
\end{table}

\subsection{The B-JEPA Architecture}

B-JEPA \citep{huang2026} is a two-stage pipeline.
\textbf{Stage~1} trains a JEPA model: a context encoder processes TF expression
tokens; a target encoder (EMA of context encoder) encodes target gene tokens;
a PoE predictor maps context representations to target representations, trained
with an ELBO objective $\mathcal{L} = \mathbb{E}[\log p(z_t \mid z_c)] - \beta
\cdot \mathrm{KL}[q(z) \| p(z)]$.
Context is formed by mean-pooling all TF representations.
\textbf{Stage~2} uses the Stage~1 predictor weights as a warm-start
$W_{\rm init} \in \mathbb{R}^{D \times G}$ for a Bayesian horseshoe regression,
trained with mean-field VI.

\subsection{Regularised Horseshoe Prior and GCV Calibration}

The regularised horseshoe \citep{piironen2017} assigns per-coefficient local
scales $\lambda_{dg} \sim C^+(0,1)$ and a global scale $\tau \sim C^+(0,\tau_0)$,
where
\begin{equation}
  \tau_0 = \frac{p_0}{D - p_0} \cdot \frac{\sigma}{\sqrt{n}},
  \label{eq:tau0}
\end{equation}
with $p_0$ the prior guess for the number of non-zero TFs per gene.
The shrinkage factor $\kappa_{dg} = 1/(1 + \lambda_{dg}^2\tau^2)$ has a
U-shaped marginal distribution under the correct posterior --- mass near 0
(signal edges) and near 1 (noise edges).
The horseshoe posterior mean is well-approximated by a GCV-calibrated ridge
estimator, which we compute analytically via SVD with no training loop required.

% ============================================================
\section{Methods}

\subsection{Experimental Setup}

All four DREAM5 networks are used throughout.
Expression matrices are standardised to zero mean and unit variance per gene.
A TF mask restricts predictors to annotated TFs; self-regulation is excluded.
All JEPA variants share Stage~1 hyperparameters: 150 epochs, Adam
($\mathrm{lr}=10^{-3}$), $\beta$-annealing from 0 to 0.1, hidden dimension 128.
Stage~2 horseshoe VI: 300 epochs (200 for exp2), Adam ($\mathrm{lr}=10^{-3}$),
$p_0 = 10$.

\subsection{JEPA Variants}

\textbf{B-JEPA}: standard pipeline as described above.
\textbf{BiJEPA}: adds a backward predictor from gene tokens to TF tokens; the
total loss is $\alpha \mathcal{L}_{\rm fwd} + (1-\alpha)\mathcal{L}_{\rm bwd}$;
$\alpha \in \{0.3, 0.5, 0.7, 0.9\}$ swept on net2.
\textbf{KAN Stage~2}: replaces the linear horseshoe regression head with
B-spline activations (grid=5, order=3, 8 basis functions per TF); uses an
identity residual initialisation to prevent warm-start misalignment.
\textbf{SM-JEPA}: stochastic masking assigns context/target tokens by biased
Bernoulli sampling with TF-bias parameter $b \in \{0.5, 0.6, 0.7, 0.8, 1.0\}$.

\subsection{Baselines and Ablations}

\textbf{GENIE3} \citep{huynh2010}: random forest feature importances, all four
networks.
\textbf{GCV-Ridge}: analytical horseshoe ridge with GCV-calibrated $\lambda$,
per-gene SVD, no training.
\textbf{Per-gene GCV}: GCV minimised independently per target gene $g$,
providing gene-specific regularisation.
\textbf{Inverted masking}: regresses TF expression on non-TF gene expression
(reverse direction); edge score is the transposed coefficient.
\textbf{Three-way ensemble}: rank-normalised weighted combination of GCV,
per-gene GCV, and B-JEPA; weights swept on a simplex grid.

\subsection{Evaluation}

Primary metric: AUPR on gold-standard edges.
All ensemble weights are optimised on the full dataset (no held-out set) ---
results represent an upper bound on ensemble performance.
For the failure diagnosis, we report Stage~1 KL at convergence, mean $|W_{\rm init}|$
across all TF--gene pairs, and the posterior shrinkage profile $p(\kappa)$
for VI vs NUTS.

% ============================================================
\section{Results}

\subsection{GCV-Ridge Matches or Beats GENIE3 on All Networks}

Table~\ref{tab:main} shows AUPR for all methods across all four networks.
GCV-ridge outperforms GENIE3 on the two biological networks with sufficient
data: net3 (\textit{E.~coli}, AUPR $0.142$ vs $0.099$, $+43\%$) and net2
(\textit{S.~aureus}) via ensemble ($0.046$ vs $0.011$, $+314\%$).
On net1 (synthetic, $n/D = 4.13$), GENIE3 outperforms standalone GCV-ridge
($0.259$ vs $0.237$) but the three-way ensemble reverses this ($0.262$
vs $0.259$, $+1\%$).
On net4, GCV + inverted masking ensemble edges ahead ($0.027$ vs $0.022$).

\begin{table}[h]
\centering
\caption{AUPR across all DREAM5 networks. Best score per network in bold.}
\label{tab:main}
\begin{tabular}{lcccc}
\toprule
Method & Net1 & Net2 & Net3 & Net4 \\
\midrule
GENIE3                      & 0.259          & 0.011 & 0.099          & 0.022 \\
GCV-Ridge                   & 0.237          & 0.018 & \textbf{0.142} & 0.024 \\
B-JEPA                      & 0.204          & 0.016 & 0.083          & 0.022 \\
Per-gene GCV                & 0.229          & 0.013 & 0.131          & 0.022 \\
Inverted masking            & 0.031          & 0.002 & 0.083          & 0.027 \\
GCV + Inverted ensemble     & 0.237          & \textbf{0.046} & 0.142 & \textbf{0.027} \\
3-way ensemble (GCV+PG+BJEPA) & \textbf{0.262} & 0.030 & 0.142 & 0.025 \\
\midrule
\textit{Best GCV-based}     & \textit{0.262} & \textit{0.046} & \textit{0.142} & \textit{0.027} \\
\bottomrule
\end{tabular}
\end{table}

The $n/D$ ratio explains the ensemble composition.
On net1 ($n/D = 4.13$), the three-way ensemble optimal weights are GCV=0.65,
B-JEPA=0.35, per-gene GCV=0.0 --- B-JEPA provides orthogonal signal when the
problem is well-conditioned.
On net2 ($n/D = 1.62$), the inverted direction dominates (GCV=0.65,
inverted=0.35) and the three-way ensemble collapses to mostly per-gene GCV
(PG=0.95) --- B-JEPA contributes noise in the underdetermined regime.
On net3/4, GCV alone is optimal and no variant improves on it substantially.

\subsection{B-JEPA Fails to Improve Over GCV-Ridge}

B-JEPA underperforms GCV-ridge on every network standalone:
net1 ($-0.033$), net2 ($-0.002$), net3 ($-0.059$), net4 ($-0.002$).
The Stage~1 KL is $1.067 \pm 0.001$ across all epochs, all networks, and all
JEPA variants (B-JEPA, BiJEPA, SM-JEPA).
This is the KL of a unit Gaussian against itself --- the approximate posterior
collapses to the prior.
$W_{\rm init}$ mean magnitudes ($0.19$--$0.24$ across networks) are
indistinguishable across all JEPA variants despite varying architecture,
data dimensionality, and training duration.

\subsection{Failure Mode I: KL Collapse Under ELBO Training on Observational Data}

The ELBO objective with mean-field Gaussian $q$ and a standard Gaussian prior
$p$ on the latent code collapses when the data distribution is multivariate
normal --- as standardised gene expression approximately is.
Under this collapse, $q(z|x) \approx p(z)$ and the representations carry no
more information than a random draw from the prior.
The KL term does not decrease because the optimal solution \emph{is} the prior:
observational co-expression data contains no signal that an unsupervised JEPA
can distinguish from confounding.

The IV interpretation clarifies why.
Observational expression data is generated by a mixture of regulatory signal
(TF $\to$ gene) and shared environmental confounding (growth condition $\to$
TF and gene jointly).
A model that receives only TF expression as context cannot separate these two
sources of variation.
The ELBO is maximised by ignoring the regulatory signal entirely --- it is
swamped by the confounding variation, which is both stronger and more
compressible in latent space.

\subsection{Failure Mode II: Dense Aggregation Cannot Represent Sparse Local Edges}

B-JEPA's context aggregation is mean-pool: $z_c = \frac{1}{D}\sum_d h_d$.
This produces a single vector $z_c \in \mathbb{R}^{d_{\rm hidden}}$ that is
\emph{identical for all target genes} $g$.
The predictor then maps $(z_c, \xi_g) \to z_{t,g}$ where $\xi_g$ is a learned
gene embedding.
All per-TF-gene specificity therefore lives in the gene embedding $\xi_g$,
not in the context representation.
The resulting $W_{\rm init} = U_c^\top \Xi$ is a product of a single shared
context embedding and per-gene embeddings --- it cannot encode the sparse
pattern of which specific TF regulates which specific gene.

KAN Stage~2 confirms this: replacing the linear horseshoe head with B-spline
activations (300 epochs) gives AUPR = $0.0164$ on net2, identical to the
linear head ($0.0164$).
The Stage~2 functional form is irrelevant because the $W_{\rm init}$ geometry
is determined entirely by Stage~1.

SM-JEPA (stochastic masking) provides a partial response.
By randomly varying which tokens are context vs target, the model is forced to
build per-token representations.
The effect is modest and regime-dependent: bias=0.6 achieves AUPR=0.0219 on
net2 ($+0.0009$ over unidirectional, best bias=0.8 on net1 $+0.0018$) ---
implicit regularisation, not causal structure recovery.
BiJEPA adds a backward predictor but the KL remains at $1.067$ regardless of
the forward weight $\alpha$; best net2 result ($\alpha=0.5$, AUPR$=0.0193$)
is below unidirectional ($0.0210$).

\subsection{Failure Mode III: Mean-Field VI Collapses the Horseshoe to Ridge}

The horseshoe prior requires a bimodal posterior $p(\kappa_{dg})$: mass near 0
for signal edges and near 1 for noise edges.
Mean-field VI with a Gaussian approximation cannot represent this bimodal
structure.
The ELBO gradient with respect to the local scale $\lambda_{dg}$ drives:
\begin{equation}
  \lambda_{dg}^* = \frac{|\mu_{dg}|}{\tau} \longrightarrow 1 \quad \forall\, d,g,
\end{equation}
because the KL term penalises deviation from the half-Cauchy prior at any
scale.
The effective prior becomes $\mathcal{N}(0, \tau^2)$ --- a ridge penalty ---
for every edge, eliminating sparsification.

GCV-ridge achieves the same result analytically in $O(n D^2)$ time with no
training loop, and is therefore strictly preferable: it avoids the VI collapse
by construction and is calibrated by GCV rather than a collapsed ELBO.

\subsection{Ensemble Analysis and the Role of Data Regime}

The ensemble results reveal that B-JEPA's representations are not \emph{useless}
--- they are useful as ensemble diversifiers precisely when the problem is
well-conditioned.
At $n/D = 4.13$ (net1), the optimal three-way ensemble assigns weight 0.35 to
B-JEPA, producing the only case where GCV-based methods beat GENIE3 on net1.
At $n/D < 2$ (net2, net4), the optimal B-JEPA weight is 0 or near-0.

The inverted masking result on net2 is the strongest practical finding:
GCV + inverted ensemble achieves AUPR=0.046 ($+314\%$ over GENIE3).
The inverted direction --- regressing TF expression on non-TF gene expression ---
captures feedback signal that the forward direction misses in the
underdetermined regime.
On net4 (\textit{S.~cerevisiae}), the inverted direction \emph{outperforms}
forward GCV standalone ($0.027$ vs $0.024$), consistent with the known
prevalence of feedback regulation in yeast.

% ============================================================
\section{Causal Identification from Observational Expression Data}

This section establishes a causal framing for GRN inference and specifies the
architectural requirements for a correct causal world model.

\subsection{The Causal Hierarchy and DREAM5}

Pearl's causal hierarchy \citep{pearl2009} distinguishes three levels:
observational ($P(Y)$), interventional ($P(Y | \mathrm{do}(X))$), and
counterfactual.
DREAM5 is observational: the expression compendium records cells under varying
growth conditions, not under targeted TF perturbations.
Every existing GRN inference method --- including GENIE3 and B-JEPA --- makes
implicit causal claims (edge $d \to g$ means TF $d$ regulates gene $g$) without
identifying these claims at the interventional level.
This is the fundamental gap that motivates the IV framework below.

\subsection{Experimental Conditions as Natural Instruments}

The DREAM5 compendium was generated under diverse experimental conditions
$E$ (temperature, nutrient availability, stress response, growth phase).
These conditions satisfy the requirements for instrumental variables \citep{angrist2009}:
\begin{enumerate}
  \item \textbf{Relevance}: $E$ affects TF expression (instruments must be
    correlated with the endogenous variable).
  \item \textbf{Exclusion}: $E$ affects target gene $g$ only \emph{through}
    TF activity, not through a direct path.
  \item \textbf{Exogeneity}: $E$ is assigned independently of unobserved
    confounders in the regulatory network.
\end{enumerate}
Under these assumptions, the causal effect of TF $d$ on gene $g$ is identified
from observational data via two-stage least squares (2SLS).
The exclusion restriction is approximate: temperature, for example, directly
affects enzyme kinetics and mRNA stability, not only TF expression.
This is an honest caveat --- the IV framing provides identification up to the
quality of the instrument, not unconditionally.

\subsection{Condition-Residualised Representations}

B-JEPA's failure under observational data (Failure Mode I) has a direct
architectural fix: partial out the shared environmental variation before
computing regulatory context.
Define:
\begin{equation}
  h_{d}^{\rm resid} = h_d - \mathrm{Proj}_{\rm cond}(h_d),
\end{equation}
where $\mathrm{Proj}_{\rm cond}$ projects onto the subspace of TF expression
variation explained by experimental conditions (estimated by PCA of TF
expression).
This is the neural analog of the 2SLS first stage: the residualised
representation $h_d^{\rm resid}$ carries only the variation in TF $d$'s
activity that is \emph{not} explained by the shared environment --- i.e., the
causal regulatory signal.

\subsection{Specification of Causal B-JEPA}

The three failure modes specify five design principles for a correct causal world
model for GRN inference:
\begin{enumerate}
  \item \textbf{Residualise conditions}: partial out environmental confounding
    before context aggregation (addresses Failure I).
  \item \textbf{Sparse local output}: use sparsemax attention
    \citep{martins2016} to compute per-gene, per-TF attention weights as edge
    scores (addresses Failure II).
  \item \textbf{Expression-space prediction}: predict in expression space, not
    latent space, to make the output directly interpretable as regulatory effect
    (addresses Failure II).
  \item \textbf{Supervised signal}: use interventional data ($P(Y|\mathrm{do}(X))$)
    or IV-identified observational signal in place of the unsupervised ELBO
    (addresses Failure I).
  \item \textbf{NUTS on sparse support}: after sparsemax identifies the
    non-zero support $\mathcal{S}_g$ per gene, run NUTS on the $|\mathcal{S}_g| \times G$
    parameter horseshoe ($|\mathcal{S}_g| \approx 15$, tractable)
    to produce calibrated credible intervals (addresses Failure III).
\end{enumerate}
This architecture --- Causal B-JEPA --- is specified here as a prior claim.
Full implementation and empirical results are the subject of subsequent work.

% ============================================================
\section{Related Work}

\paragraph{JEPA for biological sequences.}
GeneJEPA \citep{genejepa2025} is the closest prior work: it applies a
Perceiver-based JEPA with EMA teacher to single-cell RNA-seq data at scale
(Tahoe-100M, 100M profiles).
GeneJEPA predicts cellular state (cell type annotation, drug response,
perturbation response) --- it does \emph{not} infer directed regulatory edges
or produce edge-level credible intervals.
Its mean-pool readout produces a global representation, making it subject to
Failure Mode~II by construction.
JEPA-DNA \citep{assael2024} applies JEPA to DNA sequences (a different modality
from expression) and does not address GRN structure.

\paragraph{Causal and object-centric JEPA.}
C-JEPA \citep{cjepa2024} introduces object-level masking for visual scene
prediction and explicitly disclaims causal identifiability (their Remark~1),
noting that predicting in latent space does not recover interventional
distributions.
Our Failure Mode~I is the gene-expression instantiation of this same observation.

\paragraph{GRN inference baselines.}
GENIE3 \citep{huynh2010} is the standard random-forest baseline for DREAM5 and
remains competitive on synthetic and well-conditioned biological networks.
Regularised horseshoe regression with GCV calibration \citep{piironen2017,golub1979}
has not, to our knowledge, been applied as a standalone GRN inference method;
our contribution is demonstrating that it outperforms both GENIE3 and B-JEPA
analytically, with no training required.

\paragraph{Bayesian GRN inference.}
NUTS-based horseshoe posteriors \citep{gelman2013,carvalho2010} produce
calibrated credible intervals but are computationally prohibitive at full network
scale ($G > 4000$).
Our NUTS results (Appendix~\ref{app:nuts}) confirm correct horseshoe geometry
on a 20-gene showcase, providing the calibration target that VI cannot reach.

% ============================================================
\section{Discussion}

\subsection{Why Dense World Models Fail for Sparse Causal Inference}

LeCun's world model framework \citep{lecun2022} is designed for dense
prediction: predict the full state of the world in latent space from a partial
observation.
GRN inference is the opposite: identify a sparse set of local directed edges
from dense, confounded observational data.
The mismatch is structural.
Mean-pool context aggregation, the ELBO objective, and the mean-field VI
approximation are all designed for dense prediction; they fail for sparse causal
identification not because they are poorly tuned but because they are solving
the wrong problem.

This generalises beyond B-JEPA.
Any self-supervised world model applied to GRN inference will face Failure~I
(observational data is confounded) and Failure~II (dense aggregation cannot
represent sparse local structure) unless it is explicitly redesigned to address
them.
The diagnostic value of this paper is in making these failure modes precise
enough that future architectures can be evaluated against them.

\subsection{When Does GCV-Ridge Suffice?}

At $n/D > 2.4$ (net1, net3), GCV-ridge achieves the best or near-best
standalone AUPR of any method tested.
At $n/D < 2$ (net2, net4), all standalone methods degrade and ensembling is
necessary.
The practical recommendation for the field: given a new expression compendium,
compute $n/D$ first.
If $n/D > 2$, GCV-ridge is the analytically calibrated baseline that requires
no training and outperforms GENIE3.
If $n/D < 2$, ensemble GCV with the inverted direction; do not rely on
self-supervised representations.

\subsection{Limitations}

Results are limited to the DREAM5 benchmark; generalisation to single-cell or
multi-condition compendium data is not demonstrated.
The inverted masking ensemble was optimised on the full dataset without a
held-out validation set; AUPR values should be interpreted as upper bounds.
The IV assumption (exclusion restriction) is approximate and adversarial with
instrument strength --- this is an honest constraint of the observational setting,
not a claim that it is satisfied exactly.

% ============================================================
\section{Conclusion}

B-JEPA fails for GRN inference for three separable, diagnosable reasons:
KL collapse under the ELBO on observational data, dense aggregation
incompatible with sparse local causal structure, and mean-field VI collapse
of the horseshoe prior to ridge.
A GCV-calibrated ridge estimator, which violates none of these conditions, is
faster and more accurate than B-JEPA and matches or beats GENIE3 on all four
DREAM5 networks.
Ensemble combinations with inverted masking and per-gene GCV provide the best
practical performance, with the ensemble composition governed by the $n/D$ regime.
The causal identification analysis establishes that DREAM5 experimental
conditions are natural instruments and specifies the five design principles of
Causal B-JEPA --- the architecture required to produce calibrated credible
intervals over regulatory edges from observational or interventional data.

% ============================================================
\section*{References}
\bibliographystyle{abbrvnat}
\bibliography{references}

% ============================================================
\appendix

\section{Network Statistics}
\label{app:networks}

Table~\ref{tab:networks} in the main text contains full network statistics.

\section{Hyperparameter Tables}
\label{app:hyperparams}

\begin{table}[h]
\centering
\caption{Shared Stage~1 hyperparameters across all JEPA variants.}
\begin{tabular}{ll}
\toprule
Parameter & Value \\
\midrule
Epochs & 150 \\
Optimiser & Adam, lr=$10^{-3}$ \\
$\beta$-annealing & linear $0 \to 0.1$ \\
Hidden dimension & 128 \\
EMA decay & 0.996 \\
Gene embedding dim & 64 \\
\bottomrule
\end{tabular}
\end{table}

\begin{table}[h]
\centering
\caption{Stage~2 horseshoe VI hyperparameters.}
\begin{tabular}{ll}
\toprule
Parameter & Value \\
\midrule
Epochs & 300 (200 for exp2) \\
Optimiser & Adam, lr=$10^{-3}$ \\
$p_0$ & 10 \\
\bottomrule
\end{tabular}
\end{table}

\section{Full Ablation Results}
\label{app:ablations}

\begin{table}[h]
\centering
\caption{Complete AUPR results for all JEPA variants on net2 ($n/D=1.62$).}
\begin{tabular}{lc}
\toprule
Method & AUPR \\
\midrule
GCV-Ridge & 0.0175 \\
B-JEPA (unidirectional) & 0.0164 \\
BiJEPA $\alpha=0.3$ & 0.0138 \\
BiJEPA $\alpha=0.5$ & 0.0193 \\
BiJEPA $\alpha=0.7$ & 0.0184 \\
BiJEPA $\alpha=0.9$ & 0.0193 \\
SM-JEPA bias=0.5 & 0.0185 \\
SM-JEPA bias=0.6 & 0.0219 \\
SM-JEPA bias=0.7 & 0.0184 \\
SM-JEPA bias=0.8 & 0.0186 \\
SM-JEPA bias=1.0 & 0.0192 \\
KAN Stage~2 & 0.0164 \\
Linear Stage~2 & 0.0164 \\
Per-gene GCV & 0.0130 \\
Inverted masking & 0.0023 \\
GCV + Inverted ensemble & 0.0462 \\
3-way ensemble (GCV+PG+BJEPA) & 0.0299 \\
\bottomrule
\end{tabular}
\end{table}

\begin{table}[h]
\centering
\caption{SM-JEPA AUPR on net1 ($n/D=4.13$).}
\begin{tabular}{lcc}
\toprule
Method & AUPR & $\Delta$ vs Uni \\
\midrule
GCV-Ridge & 0.2371 & — \\
Unidirectional B-JEPA & 0.2250 & — \\
SM-JEPA bias=0.5 & 0.2256 & $+0.0006$ \\
SM-JEPA bias=0.6 & 0.2240 & $-0.0010$ \\
SM-JEPA bias=0.7 & 0.2241 & $-0.0009$ \\
SM-JEPA bias=0.8 & 0.2268 & $+0.0018$ \\
SM-JEPA bias=1.0 & 0.2227 & $-0.0023$ \\
\bottomrule
\end{tabular}
\end{table}

\section{NUTS Diagnostics}
\label{app:nuts}

NUTS was run on a 20-gene showcase subset of net2 using PyMC with non-centred
parameterisation (500 tuning, 500 draws, 2 chains, target accept=0.9).
All parameters achieve $\hat{R} < 1.01$ and ESS $> 100$.
The posterior $\kappa$ distribution is U-shaped (bimodal near 0 and 1),
confirming correct horseshoe geometry --- in contrast to VI which produces a
unimodal distribution near $\kappa=0.5$ (ridge-equivalent).

\end{document}
